\section{Архитектура с каким-то крутым названием класса
Colossus/Monolith и
прочая.}\label{ux430ux440ux445ux438ux442ux435ux43aux442ux443ux440ux430-ux441-ux43aux430ux43aux438ux43c-ux442ux43e-ux43aux440ux443ux442ux44bux43c-ux43dux430ux437ux432ux430ux43dux438ux435ux43c-ux43aux43bux430ux441ux441ux430-colossusmonolith-ux438-ux43fux440ux43eux447ux430ux44f.}

\begin{itemize}
\tightlist
\item
  Бастион? Не, слишком ИБП.
\item
  Парсек? Это дофига. Название выглядит фантастично. Вроде, сойдёт.
\item
  Бегемот? Не, не то.
\item
  Континент? Вроде тоже норм.
\item
  Титан? Тоже норм.
\item
  Мегалит? Вторично.
\item
  Плато -- от греческого ``рюкзак'' = σάκος πλάτης. Неплохо. Отражает
  суть хранилища, а ``плато'' - звучит. И подчёркивает, что архитектура
  крайне близка к оптимуму.
\item
  Оптимус. Неплохо.
\end{itemize}

\textbf{Вывод: мне больше всего нравится Плато.}

\subsection{Глоссарий.}\label{ux433ux43bux43eux441ux441ux430ux440ux438ux439.}

Система делится на следующие логические компоненты:

\begin{itemize}
\tightlist
\item
  Координаторы.

  \begin{itemize}
  \tightlist
  \item
    Без состояния.
  \end{itemize}
\item
  Холод.

  \begin{itemize}
  \tightlist
  \item
    RAFT-аллокатор. 64 операции при записи 2Гб в секунду. Ни о чём.
  \item
    Наборы туплов дисков.
  \end{itemize}
\item
  Жара.

  \begin{itemize}
  \tightlist
  \item
    Узлы по 3 машины в идентичной аппаратной конфигурации. Есть варианты
    на HDD, есть на NVMe + NVDIMM (оптимально для видеохостингов со
    стримами).
  \end{itemize}
\item
  WAL (распределённый WAL).

  \begin{itemize}
  \tightlist
  \item
    Пока NATS + JetStream и KV, пойдёт и сам NATS, у него есть KV с
    хорошей структурой данных: индекс в памяти, данные в LSM. Прямо
    идеально под WAL-нагрузку.
  \end{itemize}
\item
  Метаданные.

  \begin{itemize}
  \tightlist
  \item
    Records: record.id -\textgreater{} record.metadata
  \item
    Fragments: fragment.id -\textgreater{} (container.id, (crc64, size,
    compression\_type, blake) - для дедупликации, что-то ещё)
  \item
    RecordsFragmentsNo: record.id -\textgreater{}
  \item
    RecordsFragments: (seed(record.id) - для шардирования, record.id,
    номер фрагмента) -\textgreater{} fragment.id
  \item
    Containers: container.id -\textgreater{} hot\_ref \textbar{}
    cold\_ref
  \item
    Tuples: tuple.id -\textgreater{} (disk1, \ldots, diskN)
  \end{itemize}
\end{itemize}

\begin{figure}
\centering
\pandocbounded{\includesvg[keepaspectratio]{./arch.svg}}
\caption{Архитектура}
\end{figure}

\subsection{Процесс обычной загрузки
данных.}\label{ux43fux440ux43eux446ux435ux441ux441-ux43eux431ux44bux447ux43dux43eux439-ux437ux430ux433ux440ux443ux437ux43aux438-ux434ux430ux43dux43dux44bux445.}

\begin{enumerate}
\def\labelenumi{\arabic{enumi}.}
\tightlist
\item
  Запись делится на фрагменты (ограничены длиной буфера в жаре). Длины
  фрагментов кроме последнего равны друг другу.
\item
  Фрагменты сохраняются в буфере. После чего лока сразу снимается.
\item
  Координатор ищет ``свободный'' буфер в жаре и берёт на него лок. После
  чего пишет в него и лок сразу снимается.
\item
  Добавляется запись в Records, добавляются записи о загруженных
  фрагментах в Fragments, добавляются связи между записью и её
  фрагментами в RecordsFragments.
\item
  Всё это производится под WAL-ом, по определённой схеме, так что при
  обрыве связи возможна докачка.
\end{enumerate}

WritePath для заливки записи состоящей из одного фрагмента, похожим
образом строится схема заливик многофрагментных записей.

\begin{figure}
\centering
\pandocbounded{\includesvg[keepaspectratio]{./write-path-single-hot.svg}}
\caption{Запись}
\end{figure}

\subsection{Процесс
чтения.}\label{ux43fux440ux43eux446ux435ux441ux441-ux447ux442ux435ux43dux438ux44f.}

Возможные форматы чтения:

\begin{itemize}
\tightlist
\item
  Полное чтение.
\item
  Чтение от сих до сих.
\end{itemize}

Обе схемы обслуживаются одним процессом:

\begin{enumerate}
\def\labelenumi{\arabic{enumi}.}
\tightlist
\item
  Координатор ищет в Records данные по запрошенному record.id чтобы
  свериться, что запрос корректен.
\item
  Вычисляет индексы i\ldots j фрагментов в записи которые нужно
  прочесть.
\item
  Читает данные из RecordsFragments начиная с (seed(record.id),
  record.id, i) до (seed(record.id), record.id, j) и получает оттуда
  {[}{]}fragment.id.
\item
  Читает фрагменты из Fragments и находит ссылки на них (в холоде и
  жаре).
\item
  Читает фрагмент за фрагментом из указанных в данных фрагментов
  контейнеров.
\end{enumerate}

\section{Физическое устройство
системы.}\label{ux444ux438ux437ux438ux447ux435ux441ux43aux43eux435-ux443ux441ux442ux440ux43eux439ux441ux442ux432ux43e-ux441ux438ux441ux442ux435ux43cux44b.}

В данном разделе мы покажем как устроены компоненты физически.

\subsection{Жара.}\label{ux436ux430ux440ux430.}

\begin{itemize}
\tightlist
\item
  Жара состоит из узлов.
\item
  Каждый узел состоит из 3-х идентичных машин в разных ЦОД-ах.
\item
  Каждая машина держит в себе набор NVMe и плашки энергонезависимой
  памяти (есть и дешёвые варианты с батарейкой и сбросом данных на
  флешки после обесточивания).
\item
  Каждому диску \(D_i\) соответствует область фиксированного размера на
  энергонезависимой памяти: \(M_i\).
\item
  Каждый диск \(D_i\) делится на буферы фиксированного и одинакового
  размера (64Миб), \(D_i = \left( C_{i1},\ldots, C_{iM} \right)\), а
  \(M_i = \left( M_{i1}, \ldots, M_{iM}\right)\), тоже с одинаковым
  длинами \(M_{ij}\).
\end{itemize}

Что хранится в \(M_{ij}\): метаданные буфера (индекс контейнера, текущая
позиция записи, текущий объём записанных данных -- он может быть меньше
позиции записи, если где-то были неудачи) + буффер по LBA: для того
чтобы не фрагментировать данные дополняя их нулями и не перезаписывать
LBA мы храним недозаполненные LBA в памяти и сбрасываем на NVMe после
заполнения.

Запись осуществляется в один из буферов с достаточным свободным местом.
При записи берётся лок на буфер, так что больше туда никто в этот момент
писать не будет. Такой подход выбран для ``плотного'' наполнения LBA.
Количество 64Мб буферов на одном 4Тб устройстве равно 64Кб штук, этого
уже вполне достаточно, а всякие мультипуты делают такое ещё более
безболезненным.

\begin{figure}
\centering
\pandocbounded{\includesvg[keepaspectratio]{./hot.svg}}
\caption{жара}
\end{figure}

\subsubsection{Альтернатива -- всё то же самое без
NVDIMM.}\label{ux430ux43bux44cux442ux435ux440ux43dux430ux442ux438ux432ux430-ux432ux441ux451-ux442ux43e-ux436ux435-ux441ux430ux43cux43eux435-ux431ux435ux437-nvdimm.}

Придётся хранить базку (rocks, bolt и т.п.) с метаданными и смириться с
фрагментацией -- лучше дописать недозаполненный LBA нулями, чем делать
его перезаписи, это убивает ресурс SSD. Ну и банально медленнее.

\subsubsection{Альтернатива - диск как кольцевой
буфер.}\label{ux430ux43bux44cux442ux435ux440ux43dux430ux442ux438ux432ux430---ux434ux438ux441ux43a-ux43aux430ux43a-ux43aux43eux43bux44cux446ux435ux432ux43eux439-ux431ux443ux444ux435ux440.}

В этом случае мы можем использовать и HDD. Схема такая:

\begin{itemize}
\tightlist
\item
  Диск делится на буферы по 64Мб.
\item
  Мы начинаем писать в первый буфер. Если очередной фрагмент не
  помещается в остаток текущего буфера, мы запечатываем текущий и
  переходим к заполнению следующего.
\end{itemize}

\subsection{Холод.}\label{ux445ux43eux43bux43eux434.}

Узел холода состоит из RAFT-аллокатора и кортежей дисков.

Кортеж \(T = \left( D_1, \ldots, D_N \right)\)

Каждый \(D_i\), как и в случае жары, делится на блоки одинакового (между
всеми дисками кортежа) размера:

\[
D_i = \left(C_{i1}, \ldots, C_{iK}\right)
\]

Страйпом \(S_j\) в кортеже T мы назовём набор блоков одинакового индекса
на дисках кортежа:

\[
S_j = \left(C_{1j}, \ldots, C_{Nj} \right)
\]

Т.е. диски и кортежы в страйпе выглядят вот так

\begin{figure}
\centering
\pandocbounded{\includesvg[keepaspectratio]{./tuple.svg}}
\caption{Кортеж}
\end{figure}

В аллокаторе кортежу соответствует битмап, а страйпу в кортеже -- 1 бит.
1Эб дискового пространства будет соответствовать 2Гб битмап.

Таким образом, задача нахождения страйпа для переноса из горячего в
холодное хранилище будет заключаться в:

\begin{enumerate}
\def\labelenumi{\arabic{enumi}.}
\tightlist
\item
  Выбор кортежа дисков.
\item
  Выбор страйпа на диске (один из нулевых битов).
\end{enumerate}

Все необходимые данные для 1Эб дискового пространства занимают меньше
16Гб.

При скорости миграции в 2Гб/с аллокатору нужно выполнять 64
RAFT-операции.

\begin{figure}
\centering
\pandocbounded{\includesvg[keepaspectratio]{./cold.svg}}
\caption{холод}
\end{figure}

\section{Операции в
системе.}\label{ux43eux43fux435ux440ux430ux446ux438ux438-ux432-ux441ux438ux441ux442ux435ux43cux435.}

\subsection{Идентификаторы.}\label{ux438ux434ux435ux43dux442ux438ux444ux438ux43aux430ux442ux43eux440ux44b.}

Идентификаторы выдаются координаторами. Используется следующая методика:
для каждого типа идентификатора мы храним счётчик u128 на etcd:
record\_id, fragment\_id, container\_id и т.д. Координатор сразу
захватывает диапазон. Скажем, 100 000 или миллион. И затем выдаёт
значения из этого диапазона.

\subsection{Миграция из жары в
холод.}\label{ux43cux438ux433ux440ux430ux446ux438ux44f-ux438ux437-ux436ux430ux440ux44b-ux432-ux445ux43eux43bux43eux434.}

Она начинается когда узел жары обращается к координатору с задачей
перевоза.

В этом случае, координатор жары:

\begin{enumerate}
\def\labelenumi{\arabic{enumi}.}
\tightlist
\item
  Вычитывает все \textless=64Мб буфера.
\item
  Запрашивает у аллокатора страйп.
\item
  Вычисляет контрольные блоки в дополнение к имеющимся блокам данных.
\item
  Раскидывает данные по блокам страйпа.
\item
  Подтвержает занятие страйпа.
\item
  Подменяет ссылку в Containers на ColdRef - текущий страйп.
\end{enumerate}

Естественно, этот процесс чуть более хитрый, проходит под контролем
WAL-а, приходится запускать пульс в аллокатор чтобы он не сбросил
занятие страйпа, но по сути всё то же самое.

\subsection{Стриминг.}\label{ux441ux442ux440ux438ux43cux438ux43dux433.}

Первый вариант жаркого узла рассчитан на поддержку эффективной стримовой
записи. В отличии от записи одного фрагмента, стримовая запись занимает
буфер и периодически обновляет владение им, наполняя его данным до
заполнения, после чего помечает буфер запечатанным и обновляет запись в
Fragments.

Подробнее о записи в Fragments: там есть режим указывающий что фрагмент
находится в заполнении. В этом случае метаданные такого фрагмента
хранятся в горячем узле. Это нужно чтобы избежать частых обновлений
записей этого фрагмента. При заполнении или при окончании стримв
запечатка так же обновляет запись в Fragments на конечное состояние
хранимое в дескрипторе буфера.

\subsection{Уплотнение данных на
страйпах.}\label{ux443ux43fux43bux43eux442ux43dux435ux43dux438ux435-ux434ux430ux43dux43dux44bux445-ux43dux430-ux441ux442ux440ux430ux439ux43fux430ux445.}

Уплотнение страйпов заключается в персборке N страйпов в M страйпов, где
M \textless{} N. В метаданных контейнера должны хранится следующие
данные: объём незанятого места (не удалось записать фрагмент во время
наполнения или фрагмент был удалён -- место фрагмента считается
свободным) и средний размер фрагмента на диске.

Алгоритм компактизации выглядит следующим образом:

\begin{enumerate}
\def\labelenumi{\arabic{enumi}.}
\tightlist
\item
  Берём один страйп являющийся кандидатом на уплотнение.
\item
  Переносим неудалённые фрагменты на заранее выбранный страйп.
\item
  Выбираем следующий страйп и наполняем текущий приёмник его
  фрагментами.
\item
  И так далее.
\end{enumerate}

При заполнении приёмника данными из источника мы можем прийти к
следующим состояниям:

\begin{itemize}
\tightlist
\item
  Источник вычитан, но в приёмнике ещё достаточно места. Тогда
  продолжаем процесс переходя к новому источнику.
\item
  Приёмник забит, но источник ещё не вычитан. Тогда подтверждаем занятие
  приёмника и ищем следующий.
\item
  Источников больше нет. Перенос данных останавливается.
\item
  Приёмников больше нет. Но это уже общая проблема кластера, т.к.
  приёмники мы выбираем из всех туплов, а не только с данного. В этом
  случае помечаем приёмники свободными, миграция сбрасывается.
\item
  Мы достигли ограничения на количество N источников. Останавливаем
  процесс.
\end{itemize}

В процесса перевоза данных мы ведём словарь fragment.id: old\_ref
-\textgreater{} new\_ref. При завершении перевоза мы начинаем проводить
перезаписи размещений фрагментов на новые со старых.

По завершении перевоза мы помечаем все вычитанные источники свободными
(обнуляем биты в битмапах кортежей).

\subsection{Удаление
фрагментов.}\label{ux443ux434ux430ux43bux435ux43dux438ux435-ux444ux440ux430ux433ux43cux435ux43dux442ux43eux432.}

Фрагменты готовятся к удалению когда на них больше не осталось ссылок в
RecordsFragments. Процедура определения этого факта весьма тяжела если
делать напрямую:

\begin{enumerate}
\def\labelenumi{\arabic{enumi}.}
\tightlist
\item
  Хранение обратного индекса fragment.id -\textgreater{} (record.id,
  индекс фрагмента). Работать такое, скорее всего будет, но базке от
  такого точно будет нелегко. Потому что здесь O(1) от RecordsFragments,
  а это самая большая коллекция у нас.
\item
  Полное сканирование RecordsFragments. Это тяжело и ненадежно: никто не
  гарантирует доступность всех шардов во время скана базы.
\item
  Подсчёт ссылок и хранение времени доступа. Достаточно надёжно:
  refcount \textless= 0 и фрагмент не трогался уже длительное время -- с
  большой вероятностью он больше ни на что не указывает. Но такой подход
  требует перезаписей метаданных фрагментов (можно делать раз в сутки:
  видим что чисталось менее 24 часов назад -- не трогаем время
  последнего чтения).
\end{enumerate}

Каждый вариант здесь обладает своими существеннейшими недостатками.

Но есть один рабочий вариант часто применяемый для ``exascale'':

\begin{itemize}
\tightlist
\item
  Запрещаем дедупать на слишком старые фрагменты. Скажем, не старее трём
  месяцев.
\item
  Раз в какое-то время делаем снапшоты RecordsFragments и закидываем в
  clickhouse/spark/flink и прочая.
\item
  Раз в какое-то время делаем outer join-ы по fragments\_id и выбираем
  те, у которых нет связей. Если они старые то они более ни к чему не
  прикреплены и их можно удалять.
\end{itemize}
